% !PRINTED 459
\weberpage{459}

Now by (7) $\beta$ is not divisible by 3, and therefore the congruence
(5) reduces, on multiplication by $\beta$, to:

\[
\alpha+\delta \equiv 0\ (\mathrm{mod}\ 3)
\]

or, by (7) and (3), to

\weq{10}{B \equiv 0\ (\mathrm{mod}\ 3).}

If in (6) we replace $\omega$ by $\omega\pm 1$, then $B$ passes over
into $B\pm 2A$, and of the three numbers $B, B+2A, B-2A$ one and only
one is divisible by 3. We therefore assume that $B$ itself is
divisible by 3; it then follows that of the three values

\weq{11}{\gamma_2(\omega), \quad \gamma_2(\omega+1) = e^{-\frac{2\pi i}{3}}\gamma_2(\omega), \quad \gamma_2(\omega-1) = e^{\frac{2\pi i}{3}}\gamma_2(\omega)}

\gesperrt{only the first} satisfies the equation

\weq{12}{\Phi_n(u, u) = 0}

Equation (12) and

\weq{13}{u^3-j(\omega) = 0}

therefore have \gesperrt{only one} common root, and this one is rationally
expressible through $j(\omega)$. With this it is proved:

1. \gesperrt{$\gamma_2(\omega)$ is a class invariant when $\omega$ is
the root of a quadratic form whose discriminant and first coefficient
are not divisible by 3, while the middle coefficient is divisible by
3.}

If we apply this result to the case $n = 2$, we obtain first of all
from equation (5), \S\,71:

\weq{14}{\Phi_2(u, u) = u^4-2u^3-495u^2+2^4\cdot 3^3\cdot 5^3 = 0.}

The equation

\[
y^2-Dx^2 = 8
\]

is, since $D \equiv 0$ or $\equiv 1\ (\mathrm{mod}\ 4)$ must hold,
solvable only for three negative values of $D$, namely

\weq{15}{
\begin{aligned}
D &= -8, & x &= 1, & y &= 0, \\
D &= -7, & x &= 1, & y &= \pm 1, \\
D &= -4, & x &= 1, & y &= \pm 2.
\end{aligned}
}

Corresponding to these three values of $D$ one may choose the
following for equation (6):

\[
\begin{aligned}
D &= -8, & \omega^2+2 &= 0, \\
D &= -7, & \omega^2+3\omega+4 &= 0, \\
D &= -4, & \omega^2+1 &= 0
\end{aligned}
\]

or
%CONT>>
